\documentclass[12pt]{article}
\usepackage[utf8]{inputenc}
\usepackage[T1]{fontenc}
\usepackage{hyperref}
\title{Wallisian Language Resources for K-12 Curriculum Development}
\author{Compiled by Claude Sonnet 5.5 (claude-sonnet-5-5) for Andr\'as Kornai}
\date{}
\begin{document}
\maketitle

\section{Introduction}
Wallisian (East Uvean, Faka'Uvea; ISO 639-3 wls; Glottocode wall1257) is a Polynesian language of the Austronesian family, spoken on Wallis Island (\textquotesingle Uvea) in Wallis and Futuna. Ethnologue-based figures give about 10,400 speakers. Livingston (2016) says the real number is much higher, around 20,000. The searches found few resources.

\section{Primary Educational Resources}
No Wallisian primers, textbooks or school programmes were found in the searches.

\section{Linguistic and Reference Materials}
\begin{itemize}
\item Andrew Livingston (2016), \emph{East Uvean: A condensed grammar}, University of Washington MA thesis: \url{https://digital.lib.washington.edu/researchworks/items/2a9fbe02-60ce-4f08-b20d-6c04f7414ab9}
\item Endangered Languages Project entry based on Livingston (2016): \url{https://endangeredlanguages.com/lang/11041}
\item Wikipedia, Wallisian language: \url{https://en.wikipedia.org/wiki/Wallisian_language}
\item \emph{\'El\'ements de grammaire wallisienne}, a 76-page French work in the PARADISEC catalogue (unnamed author, collector Nick Thieberger): \url{https://catalog.paradisec.org.au/collections/NLA1/items/SPCmf36}
\end{itemize}
No dedicated printed dictionary was identified in the search results.

\section{Religious and Community Resources}
None were found in the searches.

\section{Digital and Audio Resources}
\begin{itemize}
\item Living Dictionaries page for Faka'uvea. I did not confirm how complete it is: \url{https://livingdictionaries.app/lea-fakauvea/about}
\end{itemize}

\section{Curriculum Development Context}
Wallisian is widely spoken, but schooling in Wallis and Futuna is mainly in French. The searches turned up no teaching materials. A curriculum would need to build on the Livingston grammar, the PARADISEC materials and local expertise. The Living Dictionaries project could serve as a vocabulary base.

\end{document}